% Imporditud: tools/import_eksperimendid.py
% Allikas: Eesti füüsikaolümpiaadi eksperimendi ülesannete
%   kogu (Rael Kalda, 2023), ülesanne Ü21, lahendus L21.
\setAuthor{EFO žürii}
\setRound{piirkonnavoor}
\setYear{2013}
\setNumber{P E2}
\setDifficulty{1}
\setTopic{vedelike mehaanika}
\setVahendid{lauatennisepall, mutter, kleeplinditükk, silindriline anum, vesi, mõõtejoonlaud. Lauatennisepalli massi võib mutri massiga võrreldes tühiseks lugeda}
\setPunktid{12}
\setAllikas{Eksperimendi kogu Ü21, lahendus lk 41}
\setLahendusseis{toores}

\prob{Mutter}
Määrake mutri mass.

\hint

\solu
Mutter tuleb kleeplindiga kinnitada lauatennise palli külge (2 p). Asetades saadud keha vette, jääb see ujuma ning selle mass on võrdne väljatõrjutud vee massiga (1 p). Väljatõrjutud vee massi saame leida veetaseme tõusu järgi (1 p). Mõõdame anuma diameetri d (1 p) ning arvutame ristlõikepindala S = $\pi$$d_2$/4 ja mõõdame veetaseme tõusu $\Delta$h (1 p), kui keha vette asetati. Keha mass m = $\rho$V = $\rho$S$\Delta$h (2 p). Sobivaks vastuseks lugeda 32 $\pm$ 2 grammi (2 p) (1 p, kui vastus on vahemikus 32 $\pm$ 4 grammi). Mutri mass on samuti 32 grammi, kuna lauatennisepalli massi loeme mutri massiga võrreldes tühiseks. Kordusmõõtmised (2 p).
\probend
